\section{Virtualios mašinos projektas}

\subsection{Virtualios mašinos samprata}
\begin{description}
\item[Virtuali mašina] - tai tarsi realios mašinos kopija. Mes
surenkame mums reikalingus komponentus, tokius kaip procesorius,
atmintis, įvedimo/išvedimo įrenginiai, suteikiame jiems paprastesnę
naudotojo sąsają. Virtuali mašina realizuoja realios mašinos
komandas paprastesniu, lengviau suprantamu būdu interpretuojant virtualios
mašinos komandas kaip realios mašinos komandas ar jų rinkiniu. Taip pat 
virtuali mašina pateikia supaprastintą atminties adresavimą.
\end{description}

\subsection{Virtualios mašinos įrangos komponentai}
\begin{description}
    \item[Procesorius] turi registrus:
    \begin{description}
        \item[IC] - 2 baitų nuorodos registras. Skirtas nurodyti vykdomos virtualios mašinos komandos adresą atmintyje.
        \item[R1] - 4 baitų bendro naudojimo registras.
        \item[R2] - 4 baitų bendro naudojimo registras.
        \item[ST] - 2 baitų rodyklė į steko segmentą atmintyje.
        \item[SP] – 2 baitų registras, kuriame saugomas steko viršūnės žodžio indeksas.
        \item[SF] - 1 baitų loginis registras, saugomi aritmetinių ir loginių operacijų suformuoti požymiai. 
    \end{description}
\end{description}

\begin{description}
    \item[Procesoriaus instrukcijos] \ 
    \begin{description}
        \item[L1xy] - atminties ląstelės, kurios adresas xy (šešioliktainis skaičius) turinio kopijavimas į registrą R1. Trumpiau: $R1:= [x*16 + y]$
        \item[L2xy] - atminties ląstelės, kurios adresas xy (šešioliktainis skaičius) turinio kopijavimas į registrą R2. Trumpiau: $R2:= [x*16 + y]$
        \item[K1xy] - registro R1 reikšmė įrašoma į atminties ląstelę, kurios adresas xy (šešioliktainis skaičius). Trumpiau $[x*16 + y]:= R1$ %K - keep
        \item[K2xy] - registro R2 reikšmė įrašoma į atminties ląstelę, kurios adresas xy (šešioliktainis skaičius). Trumpiau $[x*16 + y]:= R2$
        \item[C1xy] - į registrą R1 įrašoma į konstanta xy (šešioliktainis skaičius). Trumpiau $R1:= [x*16 + y]$
        \item[C2xy] - į registrą R2 įrašoma į konstanta xy (šešioliktainis skaičius). Trumpiau $R2:= [x*16 + y]$
        \item[A1xy] Prie registro R1 pridedama reikšmė iš atminties ląstelės, kurios adresas xy (šešioliktainis skaičius). Formuoja SF požymius. Trumpiau $R1:= R1 + [x*16 + y]$.
        \item[A2xy] Prie registro R2 pridedama reikšmė iš atminties ląstelės, kurios adresas xy (šešioliktainis skaičius). Formuoja SF požymius. Trumpiau $R2:= [x*16 + y]$.
        \item[ADD1] Prie registro R1 reikšmės pridedama registro R2 reikšmė. Formuoja SF požymius.Trumpiau $R1:= R1 + R2$.
        \item[ADD2] Prie registro R2 reikšmės pridedama registro R1 reikšmė. Formuoja SF požymius.Trumpiau $R2:= R1 + R2$.
        \item[ADxy] Sudedamos registrų R1 ir R2 reikšmės bei rezultatas išsaugojamas atminties ląstelėje, kurios adresas xy
        (šešioliktainis skaičius). Formuoja SF požymius. Trumpiau $[x*16 + y]:= R1 + R2$
        \item[S1xy] Iš registro R1 atimama reikšmė iš atminties ląstelės, kurios adresas xy (šešioliktainis skaičius). Formuoja SF požymius. Trumpiau $R1:= R1 - [x*16 + y]$.
        \item[S2xy] Iš registro R2 atimama reikšmė iš atminties ląstelės, kurios adresas xy (šešioliktainis skaičius). Formuoja SF požymius. Trumpiau $R2:= R2 - [x*16 + y]$.
        \item[SUB1] Iš registro R1 reikšmės atimama registro R2 reikšmė. Formuoja SF požymius.Trumpiau $R1:= R1 - R2$.
        \item[SUB2] Iš registro R2 reikšmės atimama registro R1 reikšmė. Formuoja SF požymius.Trumpiau $R2:= R2 - R1$.
    \end{description}
\end{description}

\begin{description}
\item[Atmintis] - virtualios mašinos atmintį sudaro 16 blokų po 16 žodžių, sunumeruotų nuo 0 iki 15.
    \begin{description}
    \item[Puslapiavimas] - realios mašinos vartotojo atmintyje saugoma lentelė, kurioje surašyti realioje atmintyje esantys blokai, priskirti virtualiai mašinai. Kiekvienas žodis (0 - 15) saugo tikrus priskirtų blokų adresus. Lentelės adresas saugomas supervizorinėje atmintyje.
    \end{description}
\end{description}

\begin{description}
    \item[Virtualios mašinos bendravimo su įvedimo/išvedimo įrenginiais mechanizmo aprašymas]
Virtualiai mašinai įvedimo ir išvedimo įrenginiai atrodo lygiai taip pat, kaip
paprasti failai, esantys išorinėje atmintyje.
\end{description}